\scanpage{39}
\mspage{38}
des conditions : $\unsure{\mathcal F|\mathscr F}$ stable dans $\mathrm{Fig}(L)$ par passage à une sous-figure, et par $\mathrm{Sup}_{\leq}$ (\unsure{ces deux-ci donc|ces deux ci-dessus} renforcent d).

Je reprends ces définitions

% One long oblique cancellation runs from the left of Soit to the foot of the sheet.
\tikz[remember picture,overlay]\coordinate (ms39DefStart);
\deleted{\uline{Définition} \quad Soit $(\mathcal F,\deleted{\unsure{\sim|\leq}}\leq,\ll)$ un atelier At1--At7 (on verra pour les axiomes ultérieurs pour \unsure{la suite|$C$ ouverts|$\mathcal C$ ouverts}).}
% slanted
\margin{\uline{Non} vide équivaut à $\varnothing_{\mathcal F}\in\mathcal F'$}

\deleted{Un \uline{sous-atelier} est une \deleted{(non vide)} partie $\mathcal F'$ \inserted{\uline{non vide}} de $\mathcal F$, fermée pour $\leq$, (i.e. telle que \unsure{t|l}\deleted{\unsure{oute|te}} sous-figure d'une figure dans $\mathcal F'$ est dans $\mathcal F'$), \inserted{et stable par $\mathrm{Sup}$ majorés dans $\mathcal F'$} On dit que $\mathcal F'$ est \inserted{un sous-atelier} \uline{spacieux} s'il est \inserted{de plus} stable par $\mathrm{Sup}$ quelconques (sans les supposer majorés dans $\mathcal F'$).}\tikz[remember picture,overlay]\coordinate (ms39DefEnd);

\deleted{Ces notions \deleted{\unsure{sont abs|sont dé}} ne dépendent que de $\leq$. Mais on munira bien sur $\mathcal F'$ des \deleted{\unsure{la|le}} \uline{deux} relations $\leq,\ll$ induites.}

\tikz[remember picture,overlay]\coordinate (ms39Soit);
\deleted{Soit}
\[
\deleted{M'=\mathcal F'\cap M}
\]
\tikz[remember picture,overlay]\coordinate (ms39Strict);
\deleted{alors on voit que les objets de $M'$ sont les objets (strictement) irréductibles \inserted{(\unsure{cf.|v.} \unsure{Ch$^{\mathrm{re}}$|Ch.} \uline{IV}, p.\,\unsure{53|52})} de $\mathcal F'$ pour $\leq$\unsure{? |;|:}. En effet, si $\deleted{\unsure{X|x}}X\in\deleted{\unsure{M'|\mathcal F}}\mathcal F'$, \deleted{et si} \inserted{\deleted{\unsure{les|les m}}\unsure{comme|même} $\mathcal F'$ est \unsure{fermé|fermée}} \deleted{$X=\mathop{\mathrm{Sup}}\limits_i\unsure{F_i|A_i}$ dans,}}
\deleted{\deleted{on a aussi $X=\mathrm{Sup}\,\unsure{F_i|E_i}$ dans $\mathcal F$ (car $\mathcal F'$ est fermé dans $\mathcal F$)} \deleted{par hyp. donc $\exists i$ avec $X=F_i$, donc $X$ irr dans $\mathcal F'$} \deleted{Inversement, si $X\in\mathcal F'$ est irr} \unsure{--|.} dans $\mathcal F$, dans $\mathcal F$, les familles $\{F_i\}$ dans $\mathcal F'$, \inserted{majorées par $X$}, sont les $\unsure{=\mathrm s|\Longleftrightarrow\mathrm s}$. Comme $\mathcal F'$ \deleted{est stable par} \inserted{$\hookrightarrow\mathcal F$ commute aux} $\mathrm{Sup}$, dire que pour une telle famille, $X=\mathrm{Sup}\,F_i$ dans $\mathcal F'$, revient au $=$. \deleted{\unsure{O|P}}Dire que pour \unsure{ttes|toutes} telles familles, $\exists i$ avec $\unsure{X_i|X}=F$, est \unsure{itou|idem} pour $\mathcal F$ ou $\mathcal F'$.}

\figure{\tikz[remember picture,overlay]{
  \draw ([xshift=-4pt,yshift=-6pt]ms39DefStart) .. controls +(-1.5pt,0pt) and +(-1pt,1pt) .. ([xshift=-4pt,yshift=-2pt]ms39DefStart |- ms39DefEnd);
  \draw ([xshift=-4pt]ms39Soit) -- ([xshift=-4pt,yshift=-12pt]ms39Soit |- ms39Strict);
  \draw ([xshift=-5.6pt]ms39Soit) -- ([xshift=-5.6pt,yshift=-11pt]ms39Soit |- ms39Strict) -- ([xshift=-4pt,yshift=-12pt]ms39Soit |- ms39Strict);
}}
\deleted{\deleted{Il faudrait vérifier}}

\deleted{Dans le cas d'un \inserted{sous-}atelier spacieux, on récupère}
